% SIAM Article Template
%\providecommand\crefalias[2]{}
\documentclass[review,hidelinks,onefignum,onetabnum]{siamart251216}
%\usepackage{cleveref} 

\usepackage{comment}

\usepackage{amsmath,amssymb,subfigure}      % for subfigures
\usepackage{booktabs,multirow,threeparttable}  %for tables
\usepackage{lipsum}
\usepackage{amsfonts}
\usepackage{graphicx}
\usepackage{epstopdf}
\usepackage{algorithmic}
\usepackage{mathrsfs} 
\usepackage{multirow}

\ifpdf
  \DeclareGraphicsExtensions{.eps,.pdf,.png,.jpg}
\else
  \DeclareGraphicsExtensions{.eps}
\fi

% Add a serial/Oxford comma by default.
\newcommand{\creflastconjunction}{, and~}

% Used for creating new theorem and remark environments
\newsiamremark{remark}{Remark}
\newsiamremark{Assumption}{Assumption}
\newsiamremark{assumption}{Assumption}
\crefname{Assumption}{Assumption}{Hypotheses}
\newsiamthm{claim}{Claim}


\usepackage{amsopn}
\DeclareMathOperator{\diag}{diag}

\newcommand{\ud}{\mathrm{d}}
\newtheorem{example}{\hspace{6mm}Example}[section]
\renewcommand{\div}{\operatorname{div}}
\graphicspath{{./figure/}{../figure/}}

\newcommand{\dx}{\,{\rm d}x}
\newcommand{\ds}{\,{\rm d}s}
\newcommand{\dd}{\,{\rm d}}
\newcommand{\Oplus}{\ensuremath{\vcenter{\hbox{\scalebox{1.5}{$\oplus$}}}}}

\newcommand{\tr}{\operatorname{tr}}
\newcommand{\alt}{\operatorname{Alt}}
\newcommand{\dev}{\operatorname{dev}}
\newcommand{\sym}{\operatorname{sym}}
\newcommand{\skw}{\operatorname{skw}}
\newcommand{\spn}{\operatorname{spn}}
\newcommand{\mspn}{\operatorname{mspn}}
\newcommand{\vspn}{\operatorname{vspn}}
\newcommand{\mskw}{\operatorname{mskw}}
\newcommand{\vskw}{\operatorname{vskw}}
\newcommand{\defm}{\operatorname{def}}

%%%% Comments of authors %%%%%
\newcommand{\LC}[1]{\textcolor{red}{#1}}
\newcommand{\XH}[1]{\textcolor{blue}{#1}}
\newcommand{\RW}[1]{\textcolor{cyan}{#1}}
\newcommand{\RZ}[1]{\textcolor{brown}{#1}}



\title{A Stabilizer Free Weak Virtual Element Method for the
Stokes Problem on Polytopal Meshes\thanks{Submitted to the editors DATE.
\funding{The first author was supported by ...}}}

% Authors: full names plus addresses.
\author{Eun-Jae Park\thanks{Department of Computational Science and Engineering, 
           Yonsei University, 
           Seoul 03722, Korea
  (\email{ejpark@yonsei.ac.kr}).}
\and Feng Wang\thanks{Jiangsu Key Laboratory for NSLSCS, 
School of Mathematical Sciences, Nanjing Normal University,
Nanjing 210023, China  
  (\email{fengwang@live.cn}).}
\and Ruishu Wang\thanks{School of Mathematics,
Jilin University, 
Changchun, Jilin 130012, China
  (\email{wangrs\_math@jlu.edu.cn}).}
}




\begin{document}



%\maketitle
%
%% REQUIRED
%\begin{abstract}
%This paper develops a novel 
%\end{abstract}
%
%% REQUIRED
%\begin{keywords}
% Error estimate
%\end{keywords}

% REQUIRED
%\begin{MSCcodes}
%65N30, 65N15, 65N12, 74B05
%\end{MSCcodes}
%65N30: Finite elements, Rayleigh-Ritz and Galerkin methods, finite methods
%65N15: Error bounds
%65N12: Stability and convergence of numerical methods
%74B05: Classical linear elasticity

\section{VEM spaces}

Let \(\Omega \subset \mathbb{R}^d\), \(d = 2, 3\) be a bounded and simply connected polytopal domain. Given an
external force field \(\boldsymbol{f} \in \boldsymbol{L}^2(\Omega)\), we consider the steady-state incompressible Stokes problem:
\begin{equation}
\left\{
\begin{aligned}
-\operatorname{div}(\boldsymbol{\varepsilon}(\boldsymbol{u})) + \nabla p &= \boldsymbol{f} \quad \text{in } \Omega, \\
\operatorname{div} \boldsymbol{u} &= 0 \quad \text{in } \Omega, \\
\boldsymbol{u} &= \boldsymbol{0} \quad \text{on } \partial\Omega.
\end{aligned}
\right.
\end{equation}


\subsection{Primal and dual meshes}
Let $\mathcal{K}_h$ be a shape-regular polytopal mesh of the domain $\Omega$. We assume each element $K\in\mathcal{K}_h$ is a star-shaped polytope with a bounded chunkiness parameter. Let $\mathcal{F}_h$ be the set of $(d-1)$-dimensional faces of $\mathcal{K}_h$, and $\mathring{\mathcal{F}}_h^{K}=\mathcal{F}_h \setminus \partial \Omega$. We call $\mathcal{K}_h$ the primal mesh.


For each $K\in\mathcal{K}_h$, let $\boldsymbol{x}_K$ be the center of the largest ball inside $K$. Connecting $\boldsymbol{x}_K$ to the vertices of $K$ gives a shape-regular refinement $\mathcal{T}_h$. Let $\mathcal{T}_h(K)$ be the set of pyramids in $K$. Let $\mathcal{F}_h^T$ be the set of $(d-1)$-dimensional faces of $\mathcal{T}_h$, and $\mathring{\mathcal{F}}_h^{T}=\mathcal{F}_h^T \setminus \partial \Omega$.

Let $F \in \partial K$ be a face with unit normal $\boldsymbol{n}_F$ and orthonormal tangent vectors $\{\boldsymbol{t}_i^{F}\}_{i=1}^{d-1}$. For any $\boldsymbol{v}\in\mathbb{R}^d$, the tangential part on $F$ is
$$
\Pi_F \boldsymbol{v}=(\boldsymbol{I}-\boldsymbol{n}_F\otimes\boldsymbol{n}_F)\boldsymbol{v}
=\sum_{i=1}^{d-1}(\boldsymbol{v}\cdot\boldsymbol{t}_i^{F})\boldsymbol{t}_i^{F}.
$$

Define the following spaces on $F$:
\begin{align*}
\mathscr T^F&:=\textrm{span}\{\boldsymbol{t}_1^F, \ldots, \boldsymbol{t}_{d-1}^F\}, \\
\mathscr T^F(\mathbb S)&:=\textrm{span}\{{\rm sym}(\boldsymbol{t}_i^F\otimes \boldsymbol{t}_j^F) \;\;\textrm{ for } 1\leq i\leq j\leq d-1\}, \\
\mathscr N^F(\mathbb S) &:= \sym(\mathscr T^F\otimes\boldsymbol{n}_F)\oplus {\rm span}\{\boldsymbol{n}_F\otimes\boldsymbol{n}_F\},\\
\mathscr T^F(\mathbb K)&:=\textrm{span}\{{\rm skw}(\boldsymbol{t}_i^F\otimes \boldsymbol{t}_j^F) \;\;\textrm{ for } 1\leq i< j\leq d-1\}.
\end{align*}
These subspaces give the following $\boldsymbol{t}$-$\boldsymbol{n}$ decompositions
\begin{align}\label{eq:decS}
\mathbb S&= \mathscr T^F(\mathbb S)\oplus \mathscr N^F(\mathbb S), \quad 
\\
\mathbb K&= \mathscr T^F(\mathbb K)\oplus \mathscr N^F(\mathbb K), \quad \mathscr N^F(\mathbb K)={\rm skw}(\mathscr T^F\otimes\boldsymbol{n}_F).
\end{align}


For \(K \in \mathcal{K}_h\), let \(\mathbb{P}_k(K)\) be the space of polynomials on \(K\) of degree at most \(k\). Let \(\mathbb{P}_k(K; \mathbb{X}) = \mathbb{P}_k(K) \otimes \mathbb{X}\), where \(\mathbb{X}\) is \(\mathbb{R}^d\) for vectors, \(\mathbb{S}\) for symmetric matrices, and \(\mathbb{K}\) for skew-symmetric matrices.

\subsection{VEM spaces}


On each element \(K\in\mathcal{K}_h\),  we introduce a local space
\begin{equation*}
\boldsymbol{V}_K = \left\{ 
\begin{aligned}
& \boldsymbol{v} \in \boldsymbol{H}(\operatorname{div}, K) \cap \boldsymbol{H}(\operatorname{curl}, K) : \boldsymbol{v} \cdot \boldsymbol{n} \in \mathbb{P}_k(F) \ \forall F \subset \partial K, \\
& \operatorname{div} \boldsymbol{v} \in \mathbb{P}_{k-1}(K), \ \operatorname{curl} \boldsymbol{v} \in \operatorname{curl} \mathbb{P}_k(K;\mathbb{R}^d)
\end{aligned}
\right\},
\end{equation*}

The global virtual element space is
\begin{equation*}
\boldsymbol{V}_h^{\operatorname{div}} := \{ \boldsymbol{v} \in \boldsymbol{H}(\operatorname{div}, \Omega), \boldsymbol{v}|_K \in \boldsymbol{V}_K \ \forall K \in \mathcal{K}_h, \boldsymbol{v} \cdot \boldsymbol{n} = 0 \text{ on } \partial \Omega \} .
\end{equation*}


The degree of freedoms (d.o.f.) for functions in \(\boldsymbol{V}_K\) are stated as
\begin{align*}
\text{Type I} \qquad & \int_F \boldsymbol{v} \cdot \boldsymbol{n} q_k \, \mathrm{d}s \qquad \forall q_k \in \mathbb{P}_k(F), F \subset \partial K, \\
\text{Type II} \qquad & \int_K \boldsymbol{v} \cdot \boldsymbol{q}_{k-2} \, \mathrm{d}\boldsymbol{x} \qquad \forall \boldsymbol{q}_{k-2} \in \mathbf{G}_{k-2}(K), \\
\text{Type III} \qquad & \int_K \boldsymbol{v} \cdot \boldsymbol{q}_k \, \mathrm{d}\boldsymbol{x} \qquad \forall \boldsymbol{q}_k \in \mathbb{P}_k(K;\mathbb{R}^d) \backslash \mathbf{G}_k(K),
\end{align*}
where \(\mathbf{G}_k(K)\) stands for the range of the gradient of the polynomial space \(\mathbb{P}_{k+1}(K)\).

On each face \(F \in \mathcal{F}_h\), we introduce a space 
$$
\boldsymbol{V}_F := \mathbb{P}_{k-1}(F;\mathscr T^F) + \mathbf{RM}(F),
$$
 where \(\mathbf{RM}(F)\) denotes the space of rigid motions, i.e., \(\mathbf{RM}(F) :=\Pi_F \mathbf{RM}(K)= \mathbb{P}_0(F; \mathscr T^F) \oplus \mathbb{P}_0(F; \mathscr T^F(\mathbb K)) \boldsymbol{x}\). The d.o.f. for functions in \(\boldsymbol{V}_F\) are
\[
\int_F \boldsymbol{v} \cdot \boldsymbol{q} \, \mathrm{d}s \quad \forall \boldsymbol{q} \in \mathbb{P}_{k-1}(F;\mathscr T^F) + \mathbf{RM}(F).
\]

Then, the element boundary space is defined as
\[
\boldsymbol{V}_h^t := \{ \boldsymbol{v}^t, \boldsymbol{v}^t|_F \in \boldsymbol{V}_F, F \in \mathcal{F}_h, \boldsymbol{v}^t = \boldsymbol{0} \text{ on } \partial \Omega \}.
\]

Combining the above two kinds of spaces together, gives our weak virtual finite element space:
\[
\boldsymbol{V}_h := \left\{ \{ \boldsymbol{v}_h^{\operatorname{div}}, \boldsymbol{v}_h^t \}, \boldsymbol{v}_h^{\operatorname{div}} \in \boldsymbol{V}_h^{\operatorname{div}}, \boldsymbol{v}_h^t \in \boldsymbol{V}_h^t \right\}.
\]

\subsection{Projection operator}

On each \(K\in\mathcal{K}_h\), we shall define a \(\boldsymbol{L}^2\) projection. For each \(\boldsymbol{v} \in \boldsymbol{V}_K\), let \(\boldsymbol{\Pi}_K^o \boldsymbol{v} \in \mathbb{P}_k(K;\mathbb{R}^d)\) satisfy
\[
\int_K \boldsymbol{\Pi}_K^o \boldsymbol{v} \cdot \boldsymbol{q} := \int_K \boldsymbol{v} \cdot \boldsymbol{q} \quad \forall \boldsymbol{q} \in \mathbb{P}_k(K;\mathbb{R}^d).
\]

Although the function in \(\boldsymbol{V}_K\) is virtual, the \(\boldsymbol{L}^2\) projection is computable. In view of Type III d.o.f., we only need to determine \(\int_K \boldsymbol{v} \cdot \nabla p_{k+1}\) for \(p_{k+1} \in \mathbb{P}_{k+1}(K)\) which is again computable through integration by parts
\[
\int_K \boldsymbol{v} \cdot \nabla p_{k+1} \, \mathrm{d}\boldsymbol{x} = - \int_K \operatorname{div} \boldsymbol{v} \, p_{k+1} \, \mathrm{d}\boldsymbol{x} + \int_{\partial K} \boldsymbol{v} \cdot \boldsymbol{n} \, p_{k+1} \, \mathrm{d}s.
\]

Note that although \(p_{k+1} \in \mathbb{P}_{k+1}(K)\), the polynomials \(\operatorname{div} \boldsymbol{v}\) and \(\boldsymbol{v} \cdot \boldsymbol{n}\) have been computed and considered as known.



\section{Weak strain}

\subsection{Weak strain space}
For $T\in\mathcal{T}_h$, the local weak strain space is defined as
\begin{align*}
&\Sigma_k(T)= \sym(\boldsymbol{n}_F\otimes V_{k}(F)) + \mathbb{P}_{k-1}(T;\mathbb S)\\
&= (\boldsymbol{n}_F\otimes\boldsymbol{n}_F\mathbb{P}_k(T)) \oplus \sym(\boldsymbol{n}_F\otimes (\mathbb{P}_{k-1}(T;\mathscr T^F) + \mathbf{RM}(F)))\oplus \mathbb{P}_{k-1}(T;\mathscr T^F(\mathbb S)),
\end{align*}
where $F\in\mathcal{F}_h^{K}\cap \partial T$ and
$$V_{k}(F) := \mathbb{P}_k(F) \boldsymbol{n}_F \oplus \boldsymbol{V}_F.$$

For $k=1$, we have
\begin{align*}
\Sigma_1(T)= (\boldsymbol{n}_F\otimes\boldsymbol{n}_F\mathbb{P}_1(T)) \oplus \sym(\boldsymbol{n}_F\otimes \mathbf{RM}(F))\oplus \mathbb{P}_{0}(T;\mathscr T^F(\mathbb S)).
\end{align*}
For $k>1$, we get
\begin{align*}
\Sigma_k(T)= (\boldsymbol{n}_F\otimes\boldsymbol{n}_F\mathbb{P}_k(T)) \oplus \sym(\boldsymbol{n}_F\otimes \mathbb{P}_{k-1}(T;\mathscr T^F))\oplus \mathbb{P}_{k-1}(T;\mathscr T^F(\mathbb S)).
\end{align*}

Note that, we naturally extend the domain of polynomials in $V_{k}(F)$ from $F$ to $K$ (or to a sub-element $T\in\mathcal T_h(K)$ sharing $F$) by assuming the function is constant along lines perpendicular to $F$. 


The trace of $\Sigma_k(T)$ onto the face $F\in\mathcal{F}_h^{K}\cap \partial T$ is then given by
\begin{equation*}
\tr_F(\Sigma_k(T))=V_{k}(F). 
\end{equation*}


We define the local symmetric tensor space on $K$ as the discontinuous polynomial space
\begin{align}\label{eq:202411132}
\Sigma_h^{-1}(K)&=\Oplus_{i=1}^n\chi_{T_i}\Sigma_k(T_i)=\{\boldsymbol{\tau}\in L^2(K;\mathbb S): \boldsymbol{\tau}|_T\in \Sigma_k(T) \textrm{ for }T\in\mathcal{T}_h(K)\},
\end{align}
where $\chi_{T_i}$ is the characteristic function of $T_i$. We  define the bubble space
\begin{equation}
\begin{aligned}
\mathbb{B}_k(\div,K;\mathbb{S})=&\{\boldsymbol{\tau}\in \Sigma_h^{-1}(K): \tr\boldsymbol{\tau} :=\boldsymbol{\tau}\boldsymbol{n}|_{\partial K}=\boldsymbol{0}\}\\
=&\Oplus_{i=1}^n\lambda_{\boldsymbol{x}_K}\chi_{T_i}(\mathbb{P}_{k-1}(T_i)\boldsymbol{n}_i\otimes\boldsymbol{n}_i) \\
\oplus
&\Oplus_{i=1}^n\lambda_{\boldsymbol{x}_K}\chi_{T_i}\sym(\mathbb{P}_{k-2}(T_i; \mathscr T^F)\otimes\boldsymbol{n}_i) \\
\oplus 
&\Oplus_{i=1}^n\chi_{T_i}\mathbb{P}_{k-1}(T_i;\mathscr T^{F_i}(\mathbb S)).
\end{aligned}
\end{equation}

 \begin{assumption}\label{meshassumption}
There exist $d(d+1)/2$ symmetric tensors $b_\ell$ such that
\begin{equation}
\mathbb{S}=\text{span}\left\{ b_{\ell}, \ell=1,\ldots,d(d+1)/2\right\},\label{equ_assumption_ddpolytope}
\end{equation}
where $\mathscr T^{F_i}(\mathbb S)=\textrm{span}\{ b_i^{m}, m=1,\ldots,d(d-1)/2\}$ for $F_i\in \partial K$, and 
\begin{equation}
\left\{ b_\ell\right\} \subset \left\{ b_i^m, i=1,\ldots,n, m=1,\ldots,d(d-1)/2\right\}.\label{equ_edge_vector_property}
\end{equation}
\end{assumption}

We now formally define the DoFs for the local stress space $\Sigma_h^{-1}(K)$.

\begin{lemma}\label{lemma_DOF}
For $K\in\mathcal{K}_h$, the space $\Sigma_h^{-1}(K)$ is uniquely determined by the following DoFs:
\begin{subequations}
\label{SigmaKDoFs}
\begin{align}
\label{SigmaKDoFs1}
\int_{F}\boldsymbol{\sigma}\boldsymbol{n}\cdot \boldsymbol{q}\ds,&\quad \boldsymbol{q}\in   V_{k}(F), F\in \partial K , \\
\label{SigmaKDoFs2}
\int_K\boldsymbol{\sigma}:\boldsymbol{\tau} \dx,& \quad  \boldsymbol{\tau}\in \mathbb{B}_k(\div,K;\mathbb{S})\setminus\Oplus_{\ell=1}^{d(d+1)/2}\left (\chi_{T_{\boldsymbol{e}_\ell}} \mathbb{P}_{k-1}(T_{\boldsymbol{e}_\ell})b_\ell \right ), \\
\label{SigmaKDoFs3}
\int_K\boldsymbol{\sigma}:\boldsymbol{\tau} \dx,& \quad  \boldsymbol{\tau}\in \mathbb{P}_{k-1}(K;\mathbb{S}),
\end{align}
\end{subequations}
where $T_{\boldsymbol{e}_\ell}\in\mathcal{T}_h(K)$ is a sub-element containing the face basis $b_\ell$ chosen in \eqref{equ_edge_vector_property}.
\end{lemma}

Let
\[
\Sigma_h^{-1} = \{\boldsymbol{\sigma} \in L^2(\Omega;\mathbb{S}) : \boldsymbol{\sigma}|_T \in \Sigma_k(T), T \in \mathcal{T}_h\}.
\]

\subsection{Weak strain}
For any \(\boldsymbol{v} = \{ \boldsymbol{v}^{\operatorname{div}}, \boldsymbol{v}^t \} \in \boldsymbol{V}_h\), 
we define \(\epsilon^w_h(\boldsymbol{v}) \in \Sigma_h^{-1}\) satisfying \( \epsilon^w_h(\boldsymbol{v})|_K=\epsilon_K^w(\boldsymbol{v})\) and
{\color{red}{
\begin{equation}
\begin{aligned}\label{equ:weakstrain1}
&(\epsilon_K^w(\boldsymbol{v}), \boldsymbol{\sigma})_K \\
&:=(\epsilon(\boldsymbol{\Pi}_K^o \boldsymbol{v}^{\operatorname{div}}),  \boldsymbol{\sigma})_K 
+ \langle ( \boldsymbol{v}^{\operatorname{div}} - \boldsymbol{\Pi}_K^o \boldsymbol{v}^{\operatorname{div}}) \cdot \boldsymbol{n}, (\boldsymbol{\sigma}\boldsymbol{n}) \cdot \boldsymbol{n} \rangle_{\partial K} \\
&+ \sum_{i=1}^{d-1} \langle (\boldsymbol{v}^t - \boldsymbol{\Pi}_K^o \boldsymbol{v}^{\operatorname{div}})\cdot \boldsymbol{\tau}_i, (\boldsymbol{\sigma}\boldsymbol{n}) \cdot \boldsymbol{\tau}_i \rangle_{\partial K},
\end{aligned}
\end{equation}
}}
for all $\boldsymbol{\sigma}\in \Sigma_h^{-1}(K)$.

{\color{gray}{
Note that since $\boldsymbol{\sigma}$ is discontinuous on $K$, 
\begin{equation}
\begin{aligned}
&(\epsilon_K^w(\boldsymbol{v}), \boldsymbol{\sigma})_K\\
&\neq - (\boldsymbol{v}^{\operatorname{div}}, \operatorname{div} \boldsymbol{\sigma})_K + \langle \boldsymbol{v}^{\operatorname{div}} \cdot \boldsymbol{n}, (\boldsymbol{\sigma}\boldsymbol{n}) \cdot \boldsymbol{n} \rangle_{\partial K} + \sum_{i=1}^{d-1} \langle \boldsymbol{v}^t \cdot \boldsymbol{\tau}_i, (\boldsymbol{\sigma}\boldsymbol{n}) \cdot \boldsymbol{\tau}_i \rangle_{\partial K} \\
&\neq - (\boldsymbol{\Pi}_K^o \boldsymbol{v}^{\operatorname{div}}, \operatorname{div} \boldsymbol{\sigma})_K + \langle \boldsymbol{v}^{\operatorname{div}} \cdot \boldsymbol{n}, (\boldsymbol{\sigma}\boldsymbol{n}) \cdot \boldsymbol{n} \rangle_{\partial K} + \sum_{i=1}^{d-1} \langle \boldsymbol{v}^t \cdot \boldsymbol{\tau}_i, (\boldsymbol{\sigma}\boldsymbol{n}) \cdot \boldsymbol{\tau}_i \rangle_{\partial K}\\
\end{aligned}
\end{equation}
}}

Let
\[
\mathcal{J}_{\partial K}^{\boldsymbol{n}}(\boldsymbol{v}) := ((\boldsymbol{v}^{\operatorname{div}} - \boldsymbol{\Pi}_K^o \boldsymbol{v}^{\operatorname{div}}) \cdot \boldsymbol{n})\boldsymbol{n}, \quad \mathcal{J}_{\partial K}^{\boldsymbol{\tau}_i}(\boldsymbol{v}) := ((\boldsymbol{v}^t - \boldsymbol{\Pi}_K^o \boldsymbol{v}^{\operatorname{div}}) \cdot \boldsymbol{\tau}_i)\boldsymbol{\tau}_i.
\]

When the superscript omitted, it indicates the whole jump on the element boundary
\[
\mathcal{J}_{\partial K}(\boldsymbol{v}) := \mathcal{J}_{\partial K}^{\boldsymbol{n}}(\boldsymbol{v}) + \mathcal{J}_{\partial K}^{\boldsymbol{\tau}}(\boldsymbol{v}) \quad \text{with} \quad \mathcal{J}_{\partial K}^{\boldsymbol{\tau}}(\boldsymbol{v}) := \sum_{i=1}^{d-1} \mathcal{J}_{\partial K}^{\boldsymbol{\tau}_i}(\boldsymbol{v}).
\]

 On the face \(F \subset \partial K\), we define \((\boldsymbol{\Pi}_{\partial K}^t \boldsymbol{v})|_F \in \boldsymbol{V}_F\) as
\[
\int_F \boldsymbol{\Pi}_{\partial K}^t \boldsymbol{v} \cdot \boldsymbol{q} := \int_F \boldsymbol{v}|_K \cdot \boldsymbol{q} \quad \forall \boldsymbol{q} \in \boldsymbol{V}_F.
\]
Define the notation \(\boldsymbol{\Pi}_{\partial K}\) as
\[
\boldsymbol{\Pi}_{\partial K} \boldsymbol{v} := (\boldsymbol{v} \cdot \boldsymbol{n})\boldsymbol{n} + \boldsymbol{\Pi}_{\partial K}^t \boldsymbol{v}.
\]

With these notations, we can simply write (\ref{equ:weakstrain1}) as
\[
(\epsilon_K^w(\boldsymbol{v}), \boldsymbol{\sigma})_K = (\epsilon(\boldsymbol{\Pi}_K^o \boldsymbol{v}^{\operatorname{div}}), \boldsymbol{\sigma})_K + \langle \boldsymbol{\Pi}_{\partial K} \mathcal{J}_{\partial K}(\boldsymbol{v}), \boldsymbol{\sigma}\boldsymbol{n} \rangle_{\partial K}.
\]


For any \(\boldsymbol{v}_h \in \boldsymbol{V}_h\), we introduce a mesh dependent \(\boldsymbol{H}^1\) seminorm as
\[
\|\boldsymbol{v}_h\|_{1,h}^2 := \sum_{K \in \mathcal{K}_h} \|\boldsymbol{v}_h\|_{1,h,K}^2,
\]
where the seminorm on each element is defined as
\[
\|\boldsymbol{v}_h\|_{1,h,K}^2 := \|\epsilon (\boldsymbol{\Pi}_K^o \boldsymbol{v}_h^{\operatorname{div}})\|_K^2 + \|h^{-\frac{1}{2}}  \boldsymbol{\Pi}_{\partial K}\mathcal{J}_{\partial K}(\boldsymbol{v}_h)\|_{\partial K}^2.
\]


\begin{lemma}\label{lem:norm}
The mesh-dependent quantity
\(
\|\boldsymbol{v}_h\|_{1,h}
\)
defines a norm on the virtual element space \(\boldsymbol{V}_h\).
\end{lemma}

\begin{proof}
It is clear that \(\|\cdot\|_{1,h}\) is a seminorm. We only need to prove positive definiteness. Suppose \(\|\boldsymbol{v}_h\|_{1,h}=0\) for some \(\boldsymbol{v}_h\in\boldsymbol{V}_h\). Then for every \(K\in\mathcal{K}_h\),
\[
\epsilon (\boldsymbol{\Pi}_K^o \boldsymbol{v}_h^{\operatorname{div}})=0
\quad\text{on} \;K.
\]
and
\[
(\boldsymbol{v}^{\operatorname{div}} - \boldsymbol{\Pi}_K^o \boldsymbol{v}^{\operatorname{div}}) \cdot \boldsymbol{n}=0,
\quad\text{and}\quad
\boldsymbol{v}^t - \boldsymbol{\Pi}_{\partial K}^t \boldsymbol{\Pi}_K^o \boldsymbol{v}^{\operatorname{div}}=\mathbf{0}
\quad\text{on} \;\partial K.
\]
The first identity implies that \(\boldsymbol{\Pi}_K^o \boldsymbol{v}_h^{\operatorname{div}}=\boldsymbol{r}_K\in \mathbf{RM}(K)\) is a rigid motion vector on each element \(K\). The third identity gives, for every face \(F\subset\partial K\) and every tangential vector \(\boldsymbol{\tau}\in\mathscr{T}^F\),
\[
\boldsymbol{v}_h^t\cdot\boldsymbol{\tau}=\boldsymbol{r}_K\cdot\boldsymbol{\tau} \quad\text{on }F.
\]

Let \(F\) be an interior face shared by two elements \(K\) and \(K'\). From the above relation applied on both sides, we obtain
\[
(\boldsymbol{r}_K-\boldsymbol{r}_{K'})\cdot\boldsymbol{\tau}=0,
\quad\forall \boldsymbol{\tau}\in\mathscr{T}^F.
\]
Since \(\boldsymbol{v}_h^{\operatorname{div}}\in\boldsymbol{H}(\operatorname{div},\Omega)\), its normal component is continuous across \(F\). This implies
\[
(\boldsymbol{r}_K-\boldsymbol{r}_{K'})\cdot\boldsymbol{n}_F=0.
\]
Thus \(\boldsymbol{r}_K=\boldsymbol{r}_{K'}\) on $F$. By the connectedness of the mesh \(\mathcal{K}_h\), there exists a global rigid motion vector \(\boldsymbol{r}\) such that \(\boldsymbol{r}_K=\boldsymbol{r}\) for all \(K\in\mathcal{K}_h\).

Consequently, \(\boldsymbol{\Pi}_K^o \boldsymbol{v}_h^{\operatorname{div}}=\boldsymbol{r}\) on every \(K\). 
Since \(\boldsymbol{r}\cdot\boldsymbol{n}|_F\in \mathbb{P}_k(F)\), \(\operatorname{div}\boldsymbol{r}=0\in \mathbb{P}_{k-1}(K)\), and \(\operatorname{curl}\boldsymbol{r}\in \operatorname{curl}\mathbb{P}_k(K;\mathbb{R}^d)\), we have \(\boldsymbol{r}\in \boldsymbol{V}_K\).
The second identity gives
\[
\int_F \boldsymbol{v}^{\operatorname{div}} \cdot \boldsymbol{n} q_k \, \mathrm{d}s = \int_F \boldsymbol{r} \cdot \boldsymbol{n} q_k \, \mathrm{d}s
 \qquad \forall q_k \in \mathbb{P}_k(F), F \subset \partial K.
\]
The definition of $\boldsymbol{\Pi}_K^o$
\[
\int_K  \boldsymbol{v}^{\operatorname{div}} \cdot \boldsymbol{q} := \int_K \boldsymbol{r} \cdot \boldsymbol{q} \quad \forall \boldsymbol{q} \in \mathbb{P}_k(K;\mathbb{R}^d)
\]
yields
\begin{align*}
& \int_K \boldsymbol{v}^{\operatorname{div}} \cdot \boldsymbol{q}_{k-2} \, \mathrm{d}\boldsymbol{x} = \int_K \boldsymbol{r} \cdot \boldsymbol{q}_{k-2} \, \mathrm{d}\boldsymbol{x}
\qquad \forall \boldsymbol{q}_{k-2} \in \mathbf{G}_{k-2}(K), \\
& \int_K \boldsymbol{v}^{\operatorname{div}} \cdot \boldsymbol{q}_k \, \mathrm{d}\boldsymbol{x} =\int_K \boldsymbol{r} \cdot \boldsymbol{q}_k \, \mathrm{d}\boldsymbol{x} 
\qquad \forall \boldsymbol{q}_k \in \mathbb{P}_k(K;\mathbb{R}^d) \backslash \mathbf{G}_k(K).
\end{align*}
Using the unisolvence of the degrees of freedom of \(\boldsymbol{V}_K\), we deduce that \(\boldsymbol{v}_h^{\operatorname{div}}=\boldsymbol{r}\) on \(\Omega\).
 The boundary condition \(\boldsymbol{v}_h^{\operatorname{div}}\cdot\boldsymbol{n}=0\) on \(\partial\Omega\) then yields \(\boldsymbol{r}\cdot\boldsymbol{n}=0\) for all boundary normals \(\boldsymbol{n}\). Since the boundary normals span \(\mathbb{R}^d\), we get \(\boldsymbol{r}=\boldsymbol{0}\). Therefore \(\boldsymbol{v}_h^{\operatorname{div}}=\boldsymbol{0}\) and, from \(\boldsymbol{v}_h^t\cdot\boldsymbol{\tau}=\boldsymbol{r}\cdot\boldsymbol{\tau}=0\) for all tangential \(\boldsymbol{\tau}\), we also have \(\boldsymbol{v}_h^t=\boldsymbol{0}\). Hence \(\boldsymbol{v}_h=\boldsymbol{0}\), which proves that \(\|\cdot\|_{1,h}\) is a norm on \(\boldsymbol{V}_h\).
\end{proof}

It follows from the Korn's inequality that
\[
\|\boldsymbol{v}_h\|_{1,h,K}^2 \eqsim \|\nabla (\boldsymbol{\Pi}_K^o \boldsymbol{v}_h^{\operatorname{div}})\|_K^2 + \|h^{-\frac{1}{2}}  \boldsymbol{\Pi}_{\partial K}\mathcal{J}_{\partial K}(\boldsymbol{v}_h)\|_{\partial K}^2.
\]

\begin{lemma} \label{lem:discinf-sup}
We have the discrete inf-sup condition
\begin{equation}\label{eq:discinfsup}
\|\boldsymbol{v}\|_{1,h}\lesssim \sup_{\boldsymbol{\sigma}_h\in\Sigma_h^{-1}}
\frac{\sum_{K\in\mathcal{K}_h}(\epsilon_K^w(\boldsymbol{v}), \boldsymbol{\sigma})_K}{\|\boldsymbol{\sigma}\|_{0,h}}\quad\forall~\boldsymbol{v}\in V_h.
\end{equation}
\end{lemma}
\begin{proof}
Let \(\boldsymbol{v}\in \boldsymbol{V}_h\) be fixed. We construct a tensor \(\boldsymbol{\sigma}\in\Sigma_h^{-1}\) satisfying the following moment conditions:
\begin{enumerate}
\item[(i)] For every \(K\in\mathcal{K}_h\) and every \(\boldsymbol{\tau}\in\mathbb{P}_{k-1}(K;\mathbb{S})\),
\[
(\boldsymbol{\sigma}, \boldsymbol{\tau})_K
= (\boldsymbol{\varepsilon}(\boldsymbol{\Pi}_K^o \boldsymbol{v}), \boldsymbol{\tau})_K.
\]
\item[(ii)] For every face \(F\in\mathcal{F}_h\) and every \(\boldsymbol{q}\in V_k(F)\),
\[
\langle \boldsymbol{\sigma}\boldsymbol{n}, \boldsymbol{q} \rangle_F
= \langle h_F^{-1} \boldsymbol{\Pi}_{\partial K}\mathcal{J}_{\partial K}(\boldsymbol{v}), \boldsymbol{q} \rangle_F.
\]
\end{enumerate}
The DoFs \eqref{SigmaKDoFs2} vanish.
By the unisolvence of the degrees of freedom for \(\Sigma_h^{-1}(K)\) (Lemma~\ref{lemma_DOF}), such a \(\boldsymbol{\sigma}\) exists and is unique.

Using the definition of \(\epsilon_K^w(\boldsymbol{v})\) and the above moment conditions, we compute
\[
(\epsilon_K^w(\boldsymbol{v}), \boldsymbol{\sigma})_K
= (\boldsymbol{\varepsilon}(\boldsymbol{\Pi}_K^o \boldsymbol{v}), \boldsymbol{\sigma})_K
+ \langle \boldsymbol{\Pi}_{\partial K}\mathcal{J}_{\partial K}(\boldsymbol{v}), \boldsymbol{\sigma}\boldsymbol{n} \rangle_{\partial K}.
\]
By condition (i), the first term equals \(\|\boldsymbol{\varepsilon}(\boldsymbol{\Pi}_K^o \boldsymbol{v})\|_K^2\). By condition (ii), the boundary term equals
\[
\sum_{F\subset\partial K} h_F^{-1} \|\boldsymbol{\Pi}_{\partial K}\mathcal{J}_{\partial K}(\boldsymbol{v})\|_F^2.
\]
Summing over all \(K\in\mathcal{K}_h\), we obtain
\[
\sum_{K\in\mathcal{K}_h} (\epsilon_K^w(\boldsymbol{v}), \boldsymbol{\sigma})_K
= \sum_{K\in\mathcal{K}_h} \|\boldsymbol{\varepsilon}(\boldsymbol{\Pi}_K^o \boldsymbol{v})\|_K^2
+ \sum_{K\in\mathcal{K}_h} \sum_{F\subset\partial K} h_F^{-1} \|\boldsymbol{\Pi}_{\partial K}\mathcal{J}_{\partial K}(\boldsymbol{v})\|_F^2
= \|\boldsymbol{v}\|_{1,h}^2.
\]
On the other hand, a standard scaling argument together with the definition of the degrees of freedom yields
\[
\|\boldsymbol{\sigma}\|_{0,h} \lesssim \|\boldsymbol{v}\|_{1,h}.
\]
Therefore,
\begin{equation*}
\begin{aligned}
&\|\boldsymbol{v}\|_{1,h}^2
\lesssim  \|\boldsymbol{\sigma}\|_{0,h}
\sup_{\boldsymbol{\sigma}_h\in\Sigma_h^{-1}}
\frac{\sum_{K\in\mathcal{K}_h} (\epsilon_K^w(\boldsymbol{v}), \boldsymbol{\sigma}_h)_K}{\|\boldsymbol{\sigma}_h\|_{0,h}}\\
&\lesssim \|\boldsymbol{v}\|_{1,h}
\sup_{\boldsymbol{\sigma}_h\in\Sigma_h^{-1}}
\frac{\sum_{K\in\mathcal{K}_h} (\epsilon_K^w(\boldsymbol{v}), \boldsymbol{\sigma}_h)_K}{\|\boldsymbol{\sigma}_h\|_{0,h}}.
\end{aligned}
\end{equation*}
Dividing by \(\|\boldsymbol{v}\|_{1,h}\) (if nonzero) gives the desired discrete inf-sup condition \eqref{eq:discinfsup}. The case \(\boldsymbol{v}=\boldsymbol{0}\) is trivial.
\end{proof}


As a corollary, we have the following lemma:

\begin{lemma}
For any \(\boldsymbol{v}_h = \{\boldsymbol{v}_h^{\operatorname{div}}, \boldsymbol{v}_h^t\} \in \boldsymbol{V}_h\), there holds that
\[
\|\boldsymbol{v}_h\|_{1,h}^2 \lesssim \|\epsilon_h^w(\boldsymbol{v}_h)\|^2.
\]
\end{lemma}
\begin{proof}
By the discrete inf-sup condition established in Lemma~\ref{lem:discinf-sup}, we have
\[
\|\boldsymbol{v}_h\|_{1,h}
\lesssim
\sup_{\boldsymbol{\sigma}\in\Sigma_h^{-1}}
\frac{\sum_{K\in\mathcal{K}_h}(\epsilon_K^w(\boldsymbol{v}_h), \boldsymbol{\sigma})_K}{\|\boldsymbol{\sigma}\|_{0,h}}.
\]
For any \(\boldsymbol{\sigma}\in\Sigma_h^{-1}\), the Cauchy--Schwarz inequality gives
\begin{equation*}
\begin{aligned}
&\sum_{K\in\mathcal{K}_h}(\epsilon_K^w(\boldsymbol{v}_h), \boldsymbol{\sigma})_K
\le
\sum_{K\in\mathcal{K}_h}\|\epsilon_K^w(\boldsymbol{v}_h)\|_K \|\boldsymbol{\sigma}\|_K\\
&\le
\left(\sum_{K\in\mathcal{K}_h}\|\epsilon_K^w(\boldsymbol{v}_h)\|_K^2\right)^{1/2}
\left(\sum_{K\in\mathcal{K}_h}\|\boldsymbol{\sigma}\|_K^2\right)^{1/2}.
\end{aligned}
\end{equation*}
Since
\[
\|\epsilon_h^w(\boldsymbol{v}_h)\|^2 = \sum_{K\in\mathcal{K}_h}\|\epsilon_K^w(\boldsymbol{v}_h)\|_K^2,
\qquad
\|\boldsymbol{\sigma}\|_{0,h}^2 = \sum_{K\in\mathcal{K}_h}\|\boldsymbol{\sigma}\|_K^2,
\]
we obtain
\[
\sum_{K\in\mathcal{K}_h}(\epsilon_K^w(\boldsymbol{v}_h), \boldsymbol{\sigma})_K
\le
\|\epsilon_h^w(\boldsymbol{v}_h)\| \|\boldsymbol{\sigma}\|_{0,h}.
\]
Dividing by \(\|\boldsymbol{\sigma}\|_{0,h}\) and taking the supremum over \(\boldsymbol{\sigma}\in\Sigma_h^{-1}\), it follows that
\[
\sup_{\boldsymbol{\sigma}\in\Sigma_h^{-1}}
\frac{\sum_{K\in\mathcal{K}_h}(\epsilon_K^w(\boldsymbol{v}_h), \boldsymbol{\sigma})_K}{\|\boldsymbol{\sigma}\|_{0,h}}
\le
\|\epsilon_h^w(\boldsymbol{v}_h)\|.
\]
Combining this with the discrete inf-sup condition yields
\[
\|\boldsymbol{v}_h\|_{1,h}
\lesssim
\|\epsilon_h^w(\boldsymbol{v}_h)\|.
\]
Squaring both sides gives the desired discrete Korn inequality:
\[
\|\boldsymbol{v}_h\|_{1,h}^2 \lesssim \|\epsilon_h^w(\boldsymbol{v}_h)\|^2.
\]
This completes the proof.
\end{proof}

\section{Numerical scheme}
Let \(S_h \subset L^2_0(\Omega)\) be a discontinuous piecewise \(P_{k-1}\) element space. Then, the corresponding discrete variational formulation for problem (1) is to seek \(\boldsymbol{u}_h = \{\boldsymbol{u}_h^{\operatorname{div}}, \boldsymbol{u}_h^t\} \in \boldsymbol{V}_h\) and \(p_h \in S_h\) satisfying
\begin{equation}\tag{14}
\left\{
\begin{aligned}
a_h(\boldsymbol{u}_h, \boldsymbol{v}_h) + b(\boldsymbol{v}_h, p_h) &= (\boldsymbol{f}, \boldsymbol{\Pi}_h^o \boldsymbol{v}_h^{\operatorname{div}}) && \forall \boldsymbol{v}_h \in \boldsymbol{V}_h, \\
b(\boldsymbol{u}_h, q_h) &= 0 && \forall q_h \in S_h,
\end{aligned}
\right.
\end{equation}
where
\begin{align*}
a_h(\boldsymbol{u}_h, \boldsymbol{v}_h) &:= (\epsilon_h^w(\boldsymbol{u}_h), \epsilon_h^w(\boldsymbol{v}_h)), \\
b(\boldsymbol{v}_h, p_h) &:= -(\operatorname{div} \boldsymbol{v}_h^{\operatorname{div}}, p_h).
\end{align*}


\begin{lemma}
For every
\(
\boldsymbol{v}_h=\{\boldsymbol{v}_h^{\operatorname{div}},\boldsymbol{v}_h^t\}\in \boldsymbol{V}_h,
\)
there holds
\[
\|\epsilon_h^w(\boldsymbol{v}_h)\|^2
\lesssim
\|\boldsymbol{v}_h\|_{1,h}^2.
\]
\end{lemma}

\begin{proof}
Fix \(K\in\mathcal{K}_h\). From the definition of the weak strain, for every
\(\boldsymbol{\sigma}\in\Sigma_h^{-1}(K)\),
\[
(\epsilon_K^w(\boldsymbol{v}_h),\boldsymbol{\sigma})_K
=
(\epsilon(\Pi_K^o\boldsymbol{v}_h^{\operatorname{div}}),\boldsymbol{\sigma})_K
+
\langle
\boldsymbol{\Pi}_{\partial K}\mathcal{J}_{\partial K}(\boldsymbol{v}_h),
\boldsymbol{\sigma}\boldsymbol{n}
\rangle_{\partial K}.
\]
Taking \(\boldsymbol{\sigma}=\epsilon_K^w(\boldsymbol{v}_h)\), we obtain
\[
\|\epsilon_K^w(\boldsymbol{v}_h)\|_K^2
=
(\epsilon(\Pi_K^o\boldsymbol{v}_h^{\operatorname{div}}),\epsilon_K^w(\boldsymbol{v}_h))_K
+
\langle
\boldsymbol{\Pi}_{\partial K}\mathcal{J}_{\partial K}(\boldsymbol{v}_h),
\epsilon_K^w(\boldsymbol{v}_h)\boldsymbol{n}
\rangle_{\partial K}.
\]
By the Cauchy--Schwarz inequality,
\[
\begin{aligned}
&
\left|
(\epsilon(\Pi_K^o\boldsymbol{v}_h^{\operatorname{div}}),
\epsilon_K^w(\boldsymbol{v}_h))_K
\right| \\
&\le
\|\epsilon(\Pi_K^o\boldsymbol{v}_h^{\operatorname{div}})\|_K
\|\epsilon_K^w(\boldsymbol{v}_h)\|_K
\le
\|\epsilon(\Pi_K^o\boldsymbol{v}_h^{\operatorname{div}})\|_K
\|\epsilon_K^w(\boldsymbol{v}_h)\|_K.
\end{aligned}
\]
For the boundary term, by the Cauchy--Schwarz inequality and the trace inequality for the finite-dimensional space \(\Sigma_h^{-1}(K)\), there exists a constant \(C_T>0\), depending only on \(k\) and the shape regularity of \(\mathcal{K}_h\), such that
\[
\left|
\langle
\boldsymbol{\Pi}_{\partial K}\mathcal{J}_{\partial K}(\boldsymbol{v}_h),
\epsilon_K^w(\boldsymbol{v}_h)\boldsymbol{n}
\rangle_{\partial K}
\right|
\le
C_T
\|h^{-1/2}\Pi_{\partial K}\mathcal{J}_{\partial K}(\boldsymbol{v}_h)\|_{\partial K}
\|\epsilon_K^w(\boldsymbol{v}_h)\|_K.
\]
Combining the above estimates gives
\[
\begin{aligned}
\|\epsilon_K^w(\boldsymbol{v}_h)\|_K^2
&\le
\|\epsilon(\Pi_K^o\boldsymbol{v}_h^{\operatorname{div}})\|_K
\|\epsilon_K^w(\boldsymbol{v}_h)\|_K \\
&\quad
+
C_T
\|h^{-1/2}\Pi_{\partial K}\mathcal{J}_{\partial K}(\boldsymbol{v}_h)\|_{\partial K}
\|\epsilon_K^w(\boldsymbol{v}_h)\|_K.
\end{aligned}
\]
If \(\|\epsilon_K^w(\boldsymbol{v}_h)\|_K\neq 0\), dividing by \(\|\epsilon_K^w(\boldsymbol{v}_h)\|_K\) yields
\[
\|\epsilon_K^w(\boldsymbol{v}_h)\|_K^2
\lesssim
\|\boldsymbol{v}_h\|_{1,h,K}^2.
\]
The case \(\|\epsilon_K^w(\boldsymbol{v}_h)\|_K=0\) is trivial. 
This completes the proof.
\end{proof}

\begin{lemma}\label{lem:10}
For all \(\boldsymbol{u}_h, \boldsymbol{v}_h \in \boldsymbol{V}_h\), it holds that
\begin{align*}
a_h(\boldsymbol{u}_h, \boldsymbol{v}_h) &\lesssim \|\boldsymbol{u}_h\|_{1,h} \|\boldsymbol{v}_h\|_{1,h}, \\
a_h(\boldsymbol{u}_h, \boldsymbol{u}_h) &\gtrsim  \|\boldsymbol{u}_h\|_{1,h}^2.
\end{align*}
\end{lemma}




\begin{lemma}\label{lem:11}
For any \(\boldsymbol{v}_h = \{\boldsymbol{v}_h^{\operatorname{div}}, \boldsymbol{v}_h^t\} \in \boldsymbol{V}_h\) and \(q_h \in S_h\), it is true that
\[
b(\boldsymbol{v}_h, q_h) \lesssim \|\boldsymbol{v}_h\|_{1,h} \|q_h\|.
\]
\end{lemma}

\begin{lemma}\label{lem:12}
For any \(p_h \in S_h\), there exists a \(\boldsymbol{v}_h = \{\boldsymbol{v}_h^{\operatorname{div}}, \boldsymbol{v}_h^t\} \in \boldsymbol{V}_h\) such that
\[
\operatorname{div} \boldsymbol{v}_h^{\operatorname{div}} = p_h \quad \text{and} \quad \|\boldsymbol{v}_h\|_{1,h} \lesssim \|p_h\|.
\]
\end{lemma}

\subsection{The interpolation}
For any \(\boldsymbol{v} \in \boldsymbol{H}^1(K)\), \(\boldsymbol{I}_K^{\operatorname{div}} \boldsymbol{v} \in \boldsymbol{V}_K\) is defined by the degree of freedoms, i.e.,
\begin{align*}
\int_F \boldsymbol{I}_K^{\operatorname{div}} \boldsymbol{v} \cdot \boldsymbol{n} \, q_k &= \int_F \boldsymbol{v} \cdot \boldsymbol{n} \, q_k && \forall q_k \in P_k(F), F \subset \partial K, \\
\int_K \boldsymbol{I}_K^{\operatorname{div}} \boldsymbol{v} \cdot \boldsymbol{q}_{k-2} &= \int_K \boldsymbol{v} \cdot \boldsymbol{q}_{k-2} && \forall \boldsymbol{q}_{k-2} \in \boldsymbol{G}_{k-2}(K), \\
\int_K \boldsymbol{I}_K^{\operatorname{div}} \boldsymbol{v} \cdot \boldsymbol{q}_k &= \int_K \boldsymbol{v} \cdot \boldsymbol{q}_k && \forall \boldsymbol{q}_k \in \boldsymbol{P}_k(K) \setminus \boldsymbol{G}_k(K).
\end{align*}

\subsection{Error equations}
In this section, we assume that \((\boldsymbol{u}, p) \in (\boldsymbol{H}^{k+1} \cap \boldsymbol{H}_0^1) \times (H^k \cap L_0^2)\) is the weak solution of the Stokes problem, and define the interpolation \(\boldsymbol{u}_I = \{\boldsymbol{u}_I^{\mathrm{div}}, \boldsymbol{u}_I^t\}\) as \(\boldsymbol{u}_I^{\mathrm{div}}|_E = \boldsymbol{I}_E^{\mathrm{div}} \boldsymbol{u}, \boldsymbol{u}_I^t|_{\partial E} = \boldsymbol{\Pi}_{\partial E}^t \boldsymbol{u}, p_I = Q_h p\), where \(Q_h\) denotes the \(L^2\) projection to the space \(S_h\). We shall estimate \(\boldsymbol{u}_h - \boldsymbol{u}_I\) and \(p_h - p_I\).

Let \(\mathcal{Q}_h: L^2(\Omega;\mathbb S)\to \Sigma_h^{-1}\) denote the \(L^2\)-projection onto \(\Sigma_h^{-1}\).
For any \(\boldsymbol{v}_h \in \boldsymbol{V}_h\), we have
\begin{equation}\label{eq:erroreq}
\begin{aligned}
& a_h(\boldsymbol{u}_h - \boldsymbol{u}_I, \boldsymbol{v}_h) + b(\boldsymbol{v}_h, p_h - p_I) \\
&= (\boldsymbol{f}, \Pi_K^o \boldsymbol{v}_h^{\operatorname{div}}) 
  - (\epsilon_h^w(\boldsymbol{u}_I), \epsilon_h^w(\boldsymbol{v}_h)) 
  + (\nabla \cdot \boldsymbol{v}_h^{\operatorname{div}}, Q_h p) \\
&= (-\nabla \cdot \epsilon(\boldsymbol{u}), \Pi_K^o \boldsymbol{v}_h^{\operatorname{div}})_K 
  + (\nabla p, \Pi_K^o \boldsymbol{v}_h^{\operatorname{div}})_K \\
&\quad - (\epsilon(\Pi_K^o \boldsymbol{I}_K^{\operatorname{div}} \boldsymbol{u}), \epsilon_h^w(\boldsymbol{v}_h))_K 
  - \langle \Pi_{\partial K} \mathcal{J}_{\partial K}(\boldsymbol{u}_I), \epsilon_h^w(\boldsymbol{v}_h)\boldsymbol{n} \rangle_{\partial K} \\
&\quad + (\nabla \cdot \boldsymbol{v}_h^{\operatorname{div}}, Q_h p) \\
&= (\epsilon(\boldsymbol{u}), \epsilon(\Pi_K^o \boldsymbol{v}_h^{\operatorname{div}}))_K 
  - \langle \epsilon(\boldsymbol{u})\boldsymbol{n}, \Pi_K^o \boldsymbol{v}_h^{\operatorname{div}} \rangle_{\partial K} \\
&\quad + (\nabla p, \Pi_K^o \boldsymbol{v}_h^{\operatorname{div}})_K 
  - (\boldsymbol{v}_h^{\operatorname{div}}, \nabla Q_h p) 
  + \langle Q_h p, \boldsymbol{v}_h^{\operatorname{div}} \cdot \boldsymbol{n} \rangle_{\partial K} \\
&\quad - (\mathcal{Q}_h \epsilon(\Pi_K^o \boldsymbol{I}_K^{\operatorname{div}} \boldsymbol{u}), \epsilon(\Pi_K^o \boldsymbol{v}_h^{\operatorname{div}})) 
  - \langle \Pi_{\partial K} \mathcal{J}_{\partial K}(\boldsymbol{v}_h), \mathcal{Q}_h \epsilon(\Pi_K^o \boldsymbol{I}_K^{\operatorname{div}} \boldsymbol{u})\boldsymbol{n} \rangle_{\partial K} \\
&\quad - \langle \Pi_{\partial K} \mathcal{J}_{\partial K}(\boldsymbol{u}_I), \epsilon_h^w(\boldsymbol{v}_h)\boldsymbol{n} \rangle_{\partial K} \\
&= (\epsilon(\boldsymbol{u}) - \mathcal{Q}_h \epsilon(\Pi_K^o \boldsymbol{I}_K^{\operatorname{div}} \boldsymbol{u}), \epsilon(\Pi_K^o \boldsymbol{v}_h^{\operatorname{div}}))_K \\
&\quad + (\Pi_K^o \boldsymbol{v}_h^{\operatorname{div}}, \nabla p - \nabla Q_h p) \\
&\quad + \langle \epsilon(\boldsymbol{u})\boldsymbol{n}, \mathcal{J}_{\partial K}(\boldsymbol{v}_h) \rangle_{\partial K} 
  + \langle Q_h p - p, \boldsymbol{v}_h^{\operatorname{div}} \cdot \boldsymbol{n} \rangle_{\partial K} \\
&\quad - \langle \Pi_{\partial K} \mathcal{J}_{\partial K}(\boldsymbol{v}_h), \mathcal{Q}_h \epsilon(\Pi_K^o \boldsymbol{I}_K^{\operatorname{div}} \boldsymbol{u})\boldsymbol{n} \rangle_{\partial K} \\
&\quad - \langle \Pi_{\partial K} \mathcal{J}_{\partial K}(\boldsymbol{u}_I), \epsilon_h^w(\boldsymbol{v}_h)\boldsymbol{n} \rangle_{\partial K}\\
&= \bigl( \boldsymbol{\varepsilon}(\boldsymbol{u}) - \mathcal{Q}_h \boldsymbol{\varepsilon}(\Pi_K^o \boldsymbol{I}_K^{\operatorname{div}} \boldsymbol{u}), \boldsymbol{\varepsilon}(\Pi_K^o \boldsymbol{v}^{\operatorname{div}}) \bigr)_K \\
&\quad + \bigl\langle \mathcal{J}_{\partial K}(\boldsymbol{v}), \boldsymbol{\varepsilon}(\boldsymbol{u})\boldsymbol{n} - \mathcal{Q}_h \boldsymbol{\varepsilon}(\Pi_K^o \boldsymbol{I}_K^{\operatorname{div}} \boldsymbol{u})\boldsymbol{n} \bigr\rangle_{\partial K} \\
&\quad - \bigl( \nabla \cdot \Pi_K^o \boldsymbol{v}^{\operatorname{div}}, p - Q_h p \bigr) 
      + \bigl\langle Q_h p - p, \boldsymbol{v}^{\operatorname{div}} \cdot \boldsymbol{n} - \Pi_K^o \boldsymbol{v}^{\operatorname{div}} \cdot \boldsymbol{n} \bigr\rangle_{\partial K} \\
&\quad - \bigl\langle \Pi_{\partial K} \mathcal{J}_{\partial K}(\boldsymbol{u}_I), \epsilon_h^w(\boldsymbol{v})\boldsymbol{n} \bigr\rangle_{\partial K}
\end{aligned}
\end{equation}
Moreover, from the discrete divergence equation, for all \(q \in S_h\),
\begin{equation}\label{eq:divergenceerror}
b(\boldsymbol{u}_h - \boldsymbol{u}_I, q) = 0.
\end{equation}

We now estimate the terms on the right-hand side of \eqref{eq:erroreq}.
First,
\begin{align*}
& (\epsilon(\boldsymbol{u}) - \mathcal{Q}_h \epsilon(\Pi_K^o \boldsymbol{I}_K^{\operatorname{div}} \boldsymbol{u}), \epsilon(\Pi_K^o \boldsymbol{v}_h^{\operatorname{div}})) \\
&\le \| \epsilon(\boldsymbol{u}) - \mathcal{Q}_h \epsilon(\Pi_K^o \boldsymbol{I}_K^{\operatorname{div}} \boldsymbol{u}) \| 
   \| \epsilon(\Pi_K^o \boldsymbol{v}_h^{\operatorname{div}}) \| \\
&\le \left( \| \epsilon(\boldsymbol{u}) - \mathcal{Q}_h \epsilon(\boldsymbol{u}) \| 
   + \| \mathcal{Q}_h \epsilon(\boldsymbol{u}) - \mathcal{Q}_h \epsilon(\Pi_K^o \boldsymbol{I}_K^{\operatorname{div}} \boldsymbol{u}) \| \right) 
   \| \boldsymbol{v}_h \|_{1,h} \\
&\le \left( C h^k \|\boldsymbol{u}\|_{k+1} 
   + \| \epsilon(\boldsymbol{u} - \Pi_K^o \boldsymbol{I}_K^{\operatorname{div}} \boldsymbol{u}) \| \right) 
   \| \boldsymbol{v}_h \|_{1,h} \\
&\le \left( C h^k \|\boldsymbol{u}\|_{k+1} 
   + \| \epsilon(\Pi_K^o (\boldsymbol{u} - \boldsymbol{I}_K^{\operatorname{div}} \boldsymbol{u})) \| \right) 
   \| \boldsymbol{v}_h \|_{1,h} \\
&\le C h^k \|\boldsymbol{u}\|_{k+1} \| \boldsymbol{v}_h \|_{1,h}.
\end{align*}
Second,
\begin{align*}
& \langle \mathcal{J}_{\partial K}(\boldsymbol{v}_h), \epsilon(\boldsymbol{u})\boldsymbol{n} - \mathcal{Q}_h \epsilon(\Pi_K^o \boldsymbol{I}_K^{\operatorname{div}} \boldsymbol{u})\boldsymbol{n} \rangle_{\partial K} \\
&\le h^{-1/2} \| \mathcal{J}_{\partial K}(\boldsymbol{v}_h) \|_{\partial K} 
   \cdot h^{1/2} \| \epsilon(\boldsymbol{u}) - \mathcal{Q}_h \epsilon(\Pi_K^o \boldsymbol{I}_K^{\operatorname{div}} \boldsymbol{u}) \|_{\partial K} \\
&\le C h^k \|\boldsymbol{u}\|_{k+1} \| \boldsymbol{v}_h \|_{1,h}.
\end{align*}
Third,
\begin{align*}
& (\nabla \cdot \Pi_K^o \boldsymbol{v}_h^{\operatorname{div}}, p - Q_h p) \\
&\le \| \nabla \cdot \Pi_K^o \boldsymbol{v}_h^{\operatorname{div}} \| \| p - Q_h p \| \\
&\le C h^k \|p\|_k \| \boldsymbol{v}_h \|_{1,h}.
\end{align*}
Fourth,
\begin{align*}
& \langle Q_h p - p, \boldsymbol{v}_h^{\operatorname{div}} \cdot \boldsymbol{n} - \Pi_K^o \boldsymbol{v}_h^{\operatorname{div}} \cdot \boldsymbol{n} \rangle_{\partial K} \\
&\le \| Q_h p - p \|_{\partial K} \| (\boldsymbol{v}_h^{\operatorname{div}} - \Pi_K^o \boldsymbol{v}_h^{\operatorname{div}}) \cdot \boldsymbol{n} \|_{\partial K} \\
&\le C h^k \|p\|_k \| \boldsymbol{v}_h \|_{1,h}.
\end{align*}
Finally,
\begin{align*}
& \langle \Pi_{\partial K} \mathcal{J}_{\partial K}(\boldsymbol{u}_I), \epsilon_h^w(\boldsymbol{v}_h)\boldsymbol{n} \rangle_{\partial K} \\
&\le \| \Pi_{\partial K} \mathcal{J}_{\partial K}(\boldsymbol{u}_I) \|_{\partial K} 
   \| \epsilon_h^w(\boldsymbol{v}_h) \|_{\partial K} \\
&\le h^{-1/2} \| \Pi_{\partial K} \mathcal{J}_{\partial K}(\boldsymbol{u}_I) \|_{\partial K} 
   \| \boldsymbol{v}_h \|_{1,h},
\end{align*}
and
\begin{align*}
& h^{-1} \| \Pi_{\partial K} \mathcal{J}_{\partial K}(\boldsymbol{u}_I) \|_{\partial K}^2 \\
&\leq h^{-1} \left( \| \boldsymbol{I}_K^{\operatorname{div}} \boldsymbol{u} - \Pi_K^o \boldsymbol{I}_K^{\operatorname{div}} \boldsymbol{u} \|_{\partial K}^2 
   + \| \boldsymbol{u} - \Pi_K^o \boldsymbol{I}_K^{\operatorname{div}} \boldsymbol{u} \|_{\partial K}^2 \right) \\
&\le \left( C h^k \|\boldsymbol{u}\|_{k+1} \right)^2.
\end{align*}
Hence,
\[
\langle \Pi_{\partial K} \mathcal{J}_{\partial K}(\boldsymbol{u}_I), \epsilon_h^w(\boldsymbol{v}_h)\boldsymbol{n} \rangle_{\partial K}
\le C h^k \|\boldsymbol{u}\|_{k+1} \| \boldsymbol{v}_h \|_{1,h}.
\]
Combining the above estimates yields the desired error bound.

\subsection{Error estimates}
Let \(e_h = \boldsymbol{u}_h - \boldsymbol{u}_I\) and \(\xi_h = p_h - p_I\). 
Taking \(\boldsymbol{v}_h = e_h\) in the error equation and using the discrete divergence constraint 
\(b(e_h, q_h)=0\) for all \(q_h\in S_h\) with \(q_h=\xi_h\), we obtain
\[
a_h(e_h, e_h) = \mathcal{R}(e_h),
\]
where \(\mathcal{R}(\boldsymbol{v}_h)\) denotes the right-hand side of the error equation. 
From the coercivity of \(a_h\), we have
\[
\|e_h\|_{1,h}^2 \lesssim a_h(e_h, e_h) = \mathcal{R}(e_h).
\]
Using the estimates for the terms of \(\mathcal{R}(\boldsymbol{v}_h)\) derived above, we get
\[
|\mathcal{R}(\boldsymbol{v}_h)| \lesssim h^k \bigl(\|\boldsymbol{u}\|_{k+1} + \|p\|_k\bigr) \|\boldsymbol{v}_h\|_{1,h}.
\]
Taking \(\boldsymbol{v}_h = e_h\), it follows that
\[
\|e_h\|_{1,h} \lesssim h^k \bigl(\|\boldsymbol{u}\|_{k+1} + \|p\|_k\bigr).
\]

Next, we estimate the pressure error. By the discrete inf-sup condition, for any \(q_h\in S_h\), there exists 
\(\boldsymbol{v}_h\in \boldsymbol{V}_h\) such that
\[
b(\boldsymbol{v}_h, q_h) \gtrsim \|\boldsymbol{v}_h\|_{1,h} \|q_h\|.
\]
From the error equation, we have
\[
b(\boldsymbol{v}_h, \xi_h) = -a_h(e_h, \boldsymbol{v}_h) + \mathcal{R}(\boldsymbol{v}_h).
\]
Thus,
\[
\|\xi_h\| 
\lesssim \sup_{\boldsymbol{v}_h\in \boldsymbol{V}_h} 
\frac{-a_h(e_h, \boldsymbol{v}_h) + \mathcal{R}(\boldsymbol{v}_h)}{\|\boldsymbol{v}_h\|_{1,h}}
\lesssim \|e_h\|_{1,h} + h^k \bigl(\|\boldsymbol{u}\|_{k+1} + \|p\|_k\bigr).
\]
Combining this with the estimate for \(\|e_h\|_{1,h}\), we obtain
\[
\|\xi_h\| \lesssim h^k \bigl(\|\boldsymbol{u}\|_{k+1} + \|p\|_k\bigr).
\]

Finally, by the triangle inequality and the interpolation error estimates,
\[
\|\boldsymbol{u} - \boldsymbol{u}_h\|_{1,h} 
\le \|\boldsymbol{u} - \boldsymbol{u}_I\|_{1,h} + \|e_h\|_{1,h} 
\lesssim h^k \bigl(\|\boldsymbol{u}\|_{k+1} + \|p\|_k\bigr),
\]
and
\[
\|p - p_h\| 
\le \|p - p_I\| + \|\xi_h\| 
\lesssim h^k \bigl(\|\boldsymbol{u}\|_{k+1} + \|p\|_k\bigr).
\]
Therefore, we have the optimal error estimate
\[
\|\boldsymbol{u} - \boldsymbol{u}_h\|_{1,h} + \|p - p_h\|
\lesssim h^k \bigl(\|\boldsymbol{u}\|_{k+1} + \|p\|_k\bigr).
\]

% !TEX root = SEVEM_Stokes.tex
\section{Numerical results}\label{sec:numerics}

\AI{For $k=1$, the velocity uses $\boldsymbol V_h$ and the pressure is
piecewise constant on the primal mesh $\mathcal K_h$.
The coefficient of the symmetric strain in \eqref{eq:stokes} is one,
which corresponds to $2\nu=1$ in the usual viscosity convention.
The weak strain \eqref{equ:weakstrain1} has the equivalent form
\[
(\varepsilon_K^w(\boldsymbol v),\boldsymbol\sigma)_K
=-(\boldsymbol v^{\operatorname{div}},\operatorname{div}_{w,K}\boldsymbol\sigma)_K
+\langle\boldsymbol v^{\operatorname{div}}\cdot\boldsymbol n,
 (\boldsymbol\sigma\boldsymbol n)\cdot\boldsymbol n\rangle_{\partial K}
+\langle\boldsymbol v^t,\Pi_F(\boldsymbol\sigma\boldsymbol n)\rangle_{\partial K},
\]
where $\boldsymbol\sigma\in\Sigma_h^{-1}(K)$ with $k=1$, and
$\operatorname{div}_{w,K}\boldsymbol\sigma\in\mathbb P_1(K;\mathbb R^d)$ satisfies
\[
(\operatorname{div}_{w,K}\boldsymbol\sigma,\boldsymbol q)_K
=-(\boldsymbol\sigma,\nabla\boldsymbol q)_K
+\langle\boldsymbol\sigma\boldsymbol n,\boldsymbol q\rangle_{\partial K}
\qquad\forall\boldsymbol q\in\mathbb P_1(K;\mathbb R^d).
\]
This definition includes the internal stress jumps through integration by
parts on the cones. Its equivalence follows from the volume projection
and the symmetry of $\boldsymbol\sigma$.}

The reported \texttt{sincosdata} test is posed on the unit square.
Tables~\ref{tab:poly}--\ref{tab:rect} give the supplied results.
\AI{The energy error is $\|\boldsymbol u_I-\boldsymbol u_h\|_{A_h}
:=a_h(\boldsymbol u_I-\boldsymbol u_h,\boldsymbol u_I-\boldsymbol u_h)^{1/2}$.
The energy and total pressure errors approach first-order convergence on
these meshes. Faster convergence in some other columns is an empirical
observation; it is not proved by \eqref{eq:final-error}.}

\AI{\begin{remark}[Numerical data to confirm]
The supplied files do not specify the exact velocity, pressure, force,
boundary data, viscosity used in the computations, or the definitions of
$\|\cdot\|_{0,h}$ and the reported maximum norms.
The implemented strain space must also be checked against
\eqref{eq:strainspace}; the original numerical description used a different
index. The table values are retained pending these checks.
The rectangular meshes do not satisfy Assumption~\ref{meshassumption}:
their face-tangential tensors span only diagonal symmetric matrices.
Table~\ref{tab:rect} therefore reports a test outside the proved mesh class.
The three mesh figure files were not supplied.
\end{remark}}

% ===== Table 1 =====
\begin{table}[htbp]
  \centering \footnotesize
  \caption{Results on polygonal meshes obtained from a triangle-mesh dual.}\label{tab:poly}
  \setlength{\tabcolsep}{3pt}\resizebox{\textwidth}{!}{\begin{tabular}{cccccccc}
    \toprule
    DoFs & $h$ & $\|\boldsymbol u_I-\boldsymbol u_h\|_{A_h}$ & $\|\boldsymbol u_I-\boldsymbol u_h\|_{0,h}$ & $\|\boldsymbol u_I-\boldsymbol u_h\|_{\infty}$ & $\|p_I-p_h\|$ & $\|p-p_h\|$ & $\|p_I-p_h\|_{\text{max}}$ \\
    \midrule
    314   & 1.43e-01 & 8.767e-01 & 2.588e-02 & 7.996e-02 & 2.071e-01 & 2.454e-01 & 5.002e-01 \\
    978   & 7.78e-02 & 4.012e-01 & 9.135e-03 & 2.645e-02 & 6.612e-02 & 1.007e-01 & 1.870e-01 \\
    3362  & 4.12e-02 & 1.884e-01 & 2.712e-03 & 7.111e-03 & 1.936e-02 & 4.463e-02 & 7.481e-02 \\
    12354 & 2.13e-02 & 9.119e-02 & 7.279e-04 & 1.817e-03 & 5.869e-03 & 2.145e-02 & 3.068e-02 \\
    47234 & 1.08e-02 & 4.490e-02 & 1.875e-04 & 4.573e-04 & 1.874e-03 & 1.061e-02 & 1.335e-02 \\
    \bottomrule
  \end{tabular}}
\end{table}

% ===== Table 2 =====
\begin{table}[htbp]
  \centering\footnotesize
  \caption{Results on triangular meshes.}\label{tab:tri}
  \setlength{\tabcolsep}{3pt}\resizebox{\textwidth}{!}{\begin{tabular}{cccccccc}
    \toprule
    DoFs & $h$ & $\|\boldsymbol u_I-\boldsymbol u_h\|_{A_h}$ & $\|\boldsymbol u_I-\boldsymbol u_h\|_{0,h}$ & $\|\boldsymbol u_I-\boldsymbol u_h\|_{\infty}$ & $\|p_I-p_h\|$ & $\|p-p_h\|$ & $\|p_I-p_h\|_{\text{max}}$ \\
    \midrule
    232   & 2.50e-01 & 4.289e-01 & 1.463e-02 & 3.022e-02 & 5.357e-01 & 5.509e-01 & 8.008e-01 \\
    880   & 1.25e-01 & 1.949e-01 & 2.959e-03 & 6.035e-03 & 2.867e-01 & 2.940e-01 & 4.262e-01 \\
    3424  & 6.25e-02 & 9.410e-02 & 6.990e-04 & 1.479e-03 & 1.450e-01 & 1.486e-01 & 2.115e-01 \\
    13504 & 3.12e-02 & 4.659e-02 & 1.727e-04 & 4.138e-04 & 7.265e-02 & 7.446e-02 & 1.045e-01 \\
    53632 & 1.56e-02 & 2.323e-02 & 4.308e-05 & 1.098e-04 & 3.634e-02 & 3.725e-02 & 5.187e-02 \\
    \bottomrule
  \end{tabular}}
\end{table}

% ===== Table 3 =====
\begin{table}[htbp]
  \centering\footnotesize
  \caption{Results on rectangular meshes, outside Assumption~\ref{meshassumption}.}\label{tab:rect}
  \setlength{\tabcolsep}{3pt}\resizebox{\textwidth}{!}{\begin{tabular}{cccccccc}
    \toprule
    DoFs & $h$ & $\|\boldsymbol u_I-\boldsymbol u_h\|_{A_h}$ & $\|\boldsymbol u_I-\boldsymbol u_h\|_{0,h}$ & $\|\boldsymbol u_I-\boldsymbol u_h\|_{\infty}$ & $\|p_I-p_h\|$ & $\|p-p_h\|$ & $\|p_I-p_h\|_{\text{max}}$ \\
    \midrule
    152    & 2.50e-01 & 6.193e-01 & 4.377e-02 & 7.423e-02 & 1.557e-02 & 1.574e-01 & 2.227e-02 \\
    560    & 1.25e-01 & 2.438e-01 & 1.277e-02 & 1.582e-02 & 9.016e-03 & 8.021e-02 & 1.778e-02 \\
    2144   & 6.25e-02 & 1.083e-01 & 3.348e-03 & 4.464e-03 & 2.688e-03 & 4.011e-02 & 6.658e-03 \\
    8384   & 3.12e-02 & 5.211e-02 & 8.477e-04 & 1.148e-03 & 6.957e-04 & 2.004e-02 & 1.943e-03 \\
    33152  & 1.56e-02 & 2.579e-02 & 2.126e-04 & 2.889e-04 & 1.752e-04 & 1.002e-02 & 5.205e-04 \\
    \bottomrule
  \end{tabular}}
\end{table}

% The figure is included only when all three supplied assets are available.
\IfFileExists{fig/poly.pdf}{%
  \IfFileExists{fig/tri.pdf}{%
    \IfFileExists{fig/rect.pdf}{%
      \begin{figure}[htbp]
      \centering
      \includegraphics[width=.31\textwidth]{fig/poly.pdf}\hfill
      \includegraphics[width=.31\textwidth]{fig/tri.pdf}\hfill
      \includegraphics[width=.31\textwidth]{fig/rect.pdf}
      \caption{Polygonal, triangular, and rectangular test meshes.}\label{fig:meshes}
      \end{figure}
    }{}%
  }{}%
}{}%


\end{document} 